\documentclass[10pt, a4paper]{article}

% ── PACKAGES ──────────────────────────────────────────────────
\usepackage[utf8]{inputenc}
\usepackage[T1]{fontenc}
\usepackage{geometry}
\geometry{
  top=2.0cm,
  bottom=2.0cm,
  left=2.2cm,
  right=2.2cm
}
\usepackage{hyperref}
\hypersetup{
  colorlinks=true,
  urlcolor=blue,
  linkcolor=blue,
  citecolor=blue,
  pdftitle={Pre-Registration P7: Field-Specified
            Developmental Arrest of Neural Crest
            Migration in Sus scrofa ---
            Cross-Fostering Experiment},
  pdfauthor={Eric Robert Lawson / OrganismCore},
  pdfsubject={Pre-registered experimental protocol
              and predictions},
  pdfkeywords={domestication syndrome, neural crest,
               attractor geometry, developmental arrest,
               Sus scrofa, cross-fostering, HPA axis,
               adrenal gland, coherence gradient field,
               pre-registration, OrganismCore}
}
\usepackage{microtype}
\usepackage{parskip}
\usepackage{booktabs}
\usepackage{array}
\usepackage{tabularx}
\usepackage{longtable}
\usepackage{xcolor}
\usepackage{mdframed}
\usepackage{enumitem}
\usepackage{titlesec}
\usepackage{lmodern}
\usepackage{fancyhdr}
\usepackage{lastpage}
\usepackage{amsmath}

% ── COLOURS ───────────────────────────────────────────────────
\definecolor{headergreen}{RGB}{20,100,60}
\definecolor{rulecolor}{RGB}{20,100,60}
\definecolor{claimbox}{RGB}{20,100,60}
\definecolor{notebox}{RGB}{240,248,240}
\definecolor{warningbox}{RGB}{255,248,220}
\definecolor{predbox}{RGB}{235,248,255}
\definecolor{scorebox}{RGB}{248,240,255}
\definecolor{failbox}{RGB}{255,235,235}
\definecolor{controlbox}{RGB}{240,255,240}

% ── HEADER / FOOTER ───────────────────────────────────────────
\pagestyle{fancy}
\fancyhf{}
\renewcommand{\headrulewidth}{0.4pt}
\fancyhead[L]{\small\color{headergreen}
  OrganismCore Pre-Registration P7 v1.0}
\fancyhead[R]{\small\color{headergreen}
  Cross-Fostering / NCC Arrest --- 2026-03-18}
\fancyfoot[C]{\small Page \thepage\ of \pageref{LastPage}}

% ── SECTION FORMAT ────────────────────────────────────────────
\titleformat{\section}
  {\large\bfseries\color{headergreen}}
  {\thesection.}{0.5em}{}[
  \vspace{-0.3em}
  \textcolor{rulecolor}{\rule{\linewidth}{0.6pt}}]
\titleformat{\subsection}
  {\normalsize\bfseries\color{headergreen}}
  {\thesubsection.}{0.5em}{}
\titleformat{\subsubsection}
  {\small\bfseries}
  {\thesubsubsection.}{0.5em}{}

\setlength{\parskip}{0.4em}
\setlength{\parindent}{0pt}

% ══════════════════════════════════════════════════════════════
\begin{document}

% ── TITLE BLOCK ───────────────────────────────────────────────
\begin{center}
  {\LARGE\bfseries\color{headergreen}
  PRE-REGISTRATION P7}\\[0.4em]
  {\Large\bfseries
  Field-Specified Developmental Arrest of\\
  Neural Crest Migration in \textit{Sus scrofa}:}\\[0.3em]
  {\large\bfseries
  The Cross-Fostering Experiment}\\[0.4em]
  {\large
  Domestication Coherence Gradient Field Applied
  from Birth\\
  to F1 Offspring of Wild-Genome \textit{Sus scrofa}\\
  Predicts Morphological Shift Toward Domestic
  Phenotype\\
  Without Genetic Modification}\\[0.6em]
  \noindent\textcolor{rulecolor}
  {\rule{\linewidth}{1pt}}\\[0.4em]
  {\normalsize
  \textbf{Author:} Eric Robert Lawson\\
  \textbf{Affiliation:} OrganismCore
  (Independent Research)\\
  \textbf{ORCID:}
  \href{https://orcid.org/0009-0002-0414-6544}
  {0009-0002-0414-6544}\\
  \textbf{Contact:}
  \href{mailto:OrganismCore@proton.me}
  {OrganismCore@proton.me}\\
  \textbf{Repository:}
  \href{https://github.com/Eric-Robert-Lawson/attractor-oncology}
  {\texttt{github.com/Eric-Robert-Lawson/
  attractor-oncology}}\\
  \textbf{Pre-registration timestamp:} 2026-03-18\\
  \textbf{Linked theory paper DOI (Zenodo):}
  [to be inserted upon Zenodo confirmation]\\
  \textbf{Document version:} 1.0\\
  \textbf{License:}
  Creative Commons Attribution 4.0
  International (CC BY 4.0)}\\[0.4em]
  \noindent\textcolor{rulecolor}
  {\rule{\linewidth}{1pt}}
\end{center}

\vspace{0.5em}

% ── CLAIM BOX ─────────────────────────────────────────────────
\begin{mdframed}[
  backgroundcolor=claimbox,
  linecolor=claimbox,
  linewidth=0pt,
  innertopmargin=0.7em,
  innerbottommargin=0.7em,
  innerleftmargin=1.0em,
  innerrightmargin=1.0em
]
\color{white}
\textbf{THE PRE-REGISTERED CLAIM}

The domestication syndrome is not produced by
genetic modification of the organism.
It is produced by the domestication coherence
gradient field ($N_{\text{domestic}}$) arresting
neural crest cell migration at the
juvenile-proximate developmental stage.
The genome does not change. The field changes.
The developmental trajectory is arrested.
The phenotype reflects that arrest.

\medskip
This pre-registration specifies: the experimental
design that tests this claim, the predicted
outcome, the scoring method, the control
conditions, and the success criteria ---
all locked before the experiment is run.

\medskip
\textbf{Primary prediction (P7):}
F1 offspring of wild-caught \textit{Sus scrofa}
(confirmed wild genome, never domesticated lineage)
raised from birth under full $N_{\text{domestic}}$
via domestic foster sow will show, at 12 months,
significantly smaller adrenal glands and reduced
HPA axis reactivity compared to genotype-matched
F1 controls raised under $N_{\text{wild}}$.
Adrenal Mass Ratio (AMR) $\leq 0.80$ relative to
wild-raised controls.
Cortisol Reactivity Ratio (CRR) $\leq 0.75$
relative to wild-raised controls.
Both criteria confirmed at 12 months.
\end{mdframed}

\vspace{1em}

% ══════════════════════════════════════════════════════════════
\section{Background and Theoretical Basis}
% ══════════════════════════════════════════════════════════════

\subsection{The Standard Model and Its Failure}

The standard explanation for the domestication
syndrome --- the convergent suite of morphological
and behavioural traits (floppy ears, depigmentation,
shortened snout, smaller adrenal glands, docility,
juvenile behaviour retention) appearing in every
domesticated species --- invokes genetic change:
selective breeding shifts allele frequencies toward
domestic-type variants, which alter developmental
programs, which produce the domestic phenotype.

This explanation is incomplete for reasons that are
quantitative and structural, not merely theoretical:

\begin{enumerate}[itemsep=0.2em]
  \item Marcstr\"{o}m \& Essen-M\"{o}ller (1968)
  documented domestication syndrome traits in the
  first captive generation of wild boar
  \textit{without selective breeding}, at
  frequencies inconsistent with recessive allele
  expression in a wild founding population.

  \item The Belyaev silver fox experiment (1979)
  demonstrated that selecting only for
  HPA axis suppressibility (tameness) produces
  the full morphological syndrome including floppy
  ears, piebald coat, and shorter snout ---
  without selecting for any morphological trait.

  \item Feralization (domestic $\rightarrow$ wild)
  occurs in months to one generation.
  Re-domestication (wild $\rightarrow$ domestic)
  requires many generations of sustained selection.
  This asymmetry is not predicted by genetic models.

  \item All syndrome traits map to a single
  developmental perturbation: mild reduction in
  neural crest cell (NCC) migration (Wilkins,
  Wrangham \& Fitch, 2014). The genetic model
  cannot explain why every independent
  domestication event produces reduction in the
  same developmental program.
\end{enumerate}

\subsection{The Attractor Geometry Explanation}

The attractor geometry framework proposes:

\[
  S \;+\; N \;+\; G \;\rightarrow\; R
\]

The genome ($S$) encodes build capacity for both
wild-type and domestic-type phenotypes.
The coherence gradient field ($N$) specifies
which developmental attractor basin the genome
resolves to within the landscape geometry ($G$).
The phenotype ($R$) is the output of that resolution.

The domestication coherence gradient field
($N_{\text{domestic}}$) acts on the maternal HPA
axis, suppressing cortisol, altering the intrauterine
coherence gradient environment, and arresting
neural crest migration at the
juvenile-proximate developmental stage.
The domestic phenotype is the signature of that
arrest. The genome does not change. The field
changes. The developmental trajectory is arrested.

\textbf{This pre-registration tests the primary
prediction of this mechanism:} that applying
$N_{\text{domestic}}$ from birth to F1 offspring
of a fully wild genome will shift morphology toward
the domestic phenotype --- specifically in the
most field-responsive tissue (adrenal gland,
neural crest-derived, HPA axis endpoint) ---
without any genetic modification.

% ══════════════════════════════════════════════════════════════
\section{Experimental Design}
\label{sec:design}
% ══════════════════════════════════════════════════════════════

\subsection{Overview}

A cross-fostering experiment in which F1 offspring
of wild-caught \textit{Sus scrofa} are placed at
birth with domestic foster sows under full
$N_{\text{domestic}}$ conditions, and compared at
3, 6, and 12 months to genotype-matched F1 controls
raised under $N_{\text{wild}}$ conditions.

All animals are genotyped at birth to confirm wild
genome identity and to match genotypes across
conditions. The experiment tests whether the
coherence gradient field applied from birth ---
not genotype --- determines the degree of neural
crest developmental arrest.

\subsection{Animal Subjects}

\textbf{Source animals:}
Wild-caught \textit{Sus scrofa} (European wild boar
or feral \textit{Sus scrofa} with confirmed
multi-generational wild ancestry, minimum 3
generations without domestic management).
Source population must be confirmed by microsatellite
or SNP genotyping to have no recent admixture with
domestic breeds (admixture $< 5\%$, confirmed by
STRUCTURE or ADMIXTURE analysis).

\textbf{F1 generation:}
Offspring of wild-caught sows, sired by
wild-caught boars. Born and immediately assigned
to experimental conditions (see Section~2.3).
Minimum $n = 8$ per condition per litter cohort.
Target total: $n \geq 20$ per condition across
$\geq 3$ independent litter cohorts.

\textbf{Domestic foster sows:}
Commercial breed domestic sows (Large White or
Landrace) with confirmed domestic breed registration.
Foster sows must have successfully raised at least
one previous litter and be confirmed healthy and
lactating at time of fostering.

\subsection{Experimental Conditions}

\subsubsection{Condition A: $N_{\text{domestic}}$
from Birth (Treatment)}

F1 wild-genome piglets placed with domestic foster
sow within 6 hours of birth, before any wild maternal
imprinting.

Field components of $N_{\text{domestic}}$:
\begin{itemize}[itemsep=0.1em]
  \item Ad libitum commercial pig feed from weaning
    (21 days)
  \item Indoor housing: temperature-regulated
    (18--22\,\textdegree C), predator-free
  \item Domestic conspecific social environment
    (foster sow + domestic litter-mates)
  \item Human contact: handler interaction minimum
    30 minutes twice daily from day 1
  \item No aversive stimuli. No predation exposure.
  \item Standard veterinary care throughout
\end{itemize}

\subsubsection{Condition B: $N_{\text{wild}}$
from Birth (Control)}

F1 wild-genome piglets remain with wild-caught
biological mother in outdoor or semi-wild enclosure.

Field components of $N_{\text{wild}}$:
\begin{itemize}[itemsep=0.1em]
  \item Natural foraging supplemented to minimum
    nutritional standard (welfare requirement)
    but not ad libitum
  \item Outdoor housing with shelter available:
    ambient temperature, no climate control
  \item Wild conspecific social environment
    (biological mother + wild litter-mates)
  \item Minimal human contact (necessary
    husbandry only)
  \item Environmental complexity including substrate
    rooting, varied terrain
\end{itemize}

\subsubsection{Condition C: Domestic-Genome Control}

Domestic breed piglets (Large White or Landrace)
raised under identical $N_{\text{domestic}}$
conditions as Condition A. Provides the
domestic-genome domestic-field reference.
Confirms that measurements in Condition A that
approach Condition C values represent genuine
field-specified arrest rather than measurement
artefact. Target $n \geq 10$.

\subsection{Genotyping Protocol}

All animals genotyped within 48 hours of birth.
Minimum panel: 15 microsatellite loci standard
for Sus scrofa population genetics, plus
genome-wide SNP panel (GeneSeek Genomic
Profiler Porcine or equivalent) to confirm:
\begin{enumerate}[itemsep=0.1em]
  \item Wild ancestry ($< 5\%$ domestic admixture)
    for Condition A and B animals
  \item Genotype matching across conditions
    (Conditions A and B animals selected to have
    equivalent allele frequency distributions
    at all measured loci)
  \item No genetic outliers included in primary
    analysis
\end{enumerate}

Genotyping laboratory must be blinded to
experimental condition assignment.

% ══════════════════════════════════════════════════════════════
\section{Primary Outcome Measures}
\label{sec:outcomes}
% ══════════════════════════════════════════════════════════════

\subsection{Primary Endpoints}

The primary endpoints are neural crest-derived
tissue measurements, selected because they are
the most direct downstream indicators of NCC
migration completion and the most rapidly
responsive to field change:

\subsubsection{Adrenal Gland Mass (Primary)}

\textbf{Measurement:} Total bilateral adrenal
gland mass at necropsy (12 months), normalised
to body weight.
\[
  \text{AMR} = \frac{\bar{A}_{\text{Cond.A}}}
                    {\bar{A}_{\text{Cond.B}}}
\]
where $\bar{A}$ is mean normalised bilateral
adrenal mass within condition.

\textbf{Rationale:} The adrenal medulla is
NCC-derived (chromaffin cells originate from
NCC). Adrenal medulla mass reflects the degree
of NCC migration completion to the adrenal
anlage during embryogenesis. This is the tissue
identified by Hemmer (1990) as the earliest and
most consistent indicator of domestication-state
change in captive wild animals.

\textbf{Protocol:}
Animals humanely euthanised at 12 months.
Both adrenal glands dissected, trimmed of
periadrenal fat, weighed to nearest 0.001\,g.
Medulla and cortex dissected and weighed
separately post-fixation (4\% PFA, 48h).
Medulla mass recorded as primary measurement.
Body weight recorded at euthanasia.
Normalised adrenal medulla mass =
medulla mass / body weight $\times$ 1000
(mg/kg).

\subsubsection{HPA Axis Reactivity (Primary)}

\textbf{Measurement:} Serum cortisol response
to standardised stressor at 3, 6, and 12 months.

\textbf{Cortisol Reactivity Ratio:}
\[
  \text{CRR} = \frac{\Delta[\text{cortisol}]
  _{\text{Cond.A}}}
  {\Delta[\text{cortisol}]_{\text{Cond.B}}}
\]
where $\Delta[\text{cortisol}]$ is peak cortisol
minus baseline cortisol at T$+$30 minutes
post-stressor.

\textbf{Stressor protocol (standardised):}
Animal moved to unfamiliar pen (3m$\times$3m,
novel environment, no conspecifics visible).
Blood draw at T$=$0 (baseline, before stressor),
T$=$30 min (peak stress response),
T$=$60 min (recovery).
Serum cortisol by validated ELISA
(Enzo Life Sciences porcine cortisol kit
or equivalent with intra-assay CV $< 10\%$).

Same stressor protocol applied identically
to all conditions at identical ages.
Handlers blinded to condition assignment
during blood draws.

\subsection{Secondary Endpoints}

Secondary endpoints are measured at 3, 6,
and 12 months and reported as exploratory.
They are not included in the primary
success/failure criteria but are reported
as confirmatory evidence of the mechanism.

\subsubsection{Coat Pigmentation Pattern}

\textbf{Measurement:} Percentage of total
body surface area showing depigmentation
(white or pale patches) at 6 and 12 months.
Photographed under standardised lighting
(dorsal, lateral left, lateral right views).
Image analysis by blinded scorer using
ImageJ colour thresholding.
\[
  \text{Pigmentation Score} =
  \frac{\text{White patch area (cm}^2\text{)}}
  {\text{Total body surface area (cm}^2\text{)}}
  \times 100
\]

\subsubsection{Ear Cartilage Position}

\textbf{Measurement:} Ear position angle from
vertical plane at 6 and 12 months.
Photographed from anterior aspect, standardised
restraint. Angle measured from ear base to ear
tip relative to vertical (0\textdegree{} =
fully erect; 90\textdegree{} = fully lateral).
Mean of left and right ear.

\subsubsection{Snout Morphology Index}

\textbf{Measurement:} Snout length (nasal tip
to orbital midpoint) divided by snout width
(widest point at base of snout) at 6 and 12
months. Measured from standardised lateral
and dorsal photographs using ImageJ.
Higher ratio = more elongated (wild-type).
Lower ratio = more abbreviated (domestic-type).

\subsubsection{Hair Coarseness Score}

\textbf{Measurement:} Tactile scoring by
blinded handler (1--5 scale: 1 = fine domestic
coat, 5 = coarse wild-type bristle) at dorsal
mid-back location. Minimum 3 independent
scorers. Mean score recorded. Inter-rater
reliability confirmed (Cronbach's $\alpha
> 0.80$ required before scoring begins).

\subsubsection{Body Composition}

\textbf{Measurement:} Ultrasound backfat depth
at P2 position (6.5\,cm lateral to dorsal
midline at last rib) at 6 and 12 months.
B-mode ultrasound. Higher backfat = more
domestic metabolic program.

% ══════════════════════════════════════════════════════════════
\section{The Two Models and Their Diverging
Predictions}
\label{sec:models}
% ══════════════════════════════════════════════════════════════

This experiment is designed to distinguish
two models that make identical predictions for
behaviour but diverging predictions for morphology:

\begin{center}
\small
\begin{tabular}{p{4.5cm}p{4.5cm}p{4.5cm}}
\toprule
\textbf{Outcome} &
\textbf{Genetic model} &
\textbf{Field-arrest model} \\
\midrule
Adrenal mass in Cond.\ A vs.\ Cond.\ B
(identical genomes) &
No difference &
\textbf{Smaller in Cond.\ A}
(AMR $\leq 0.80$) \\[0.4em]
HPA reactivity in Cond.\ A vs.\ Cond.\ B &
No difference &
\textbf{Reduced in Cond.\ A}
(CRR $\leq 0.75$) \\[0.4em]
Coat depigmentation in Cond.\ A vs.\ Cond.\ B &
No difference &
\textbf{Higher in Cond.\ A} \\[0.4em]
Ear position in Cond.\ A vs.\ Cond.\ B &
No difference &
\textbf{Greater droop in Cond.\ A} \\[0.4em]
Behaviour (tameness) in Cond.\ A
vs.\ Cond.\ B &
\textbf{More tame in Cond.\ A}
(learning/imprinting) &
\textbf{More tame in Cond.\ A}
(basin occupation) \\
\bottomrule
\end{tabular}
\end{center}

\vspace{0.3em}
\textbf{The models agree on behaviour.}
Behaviour change with domestic field application
from birth is explained by learning and imprinting
under the genetic model and by basin occupation
under the field-arrest model.

\textbf{The models diverge on morphology.}
The genetic model predicts no morphological
difference between genotype-matched animals in
different fields. The field-arrest model predicts
specific morphological shifts in NCC-derived
tissues in the direction of the applied field.

The primary endpoints (adrenal mass, HPA
reactivity) are the measurements that distinguish
the models. Behaviour is not the distinguishing
test. Morphology is.

% ══════════════════════════════════════════════════════════════
\section{Success Criteria}
\label{sec:success}
% ══════════════════════════════════════════════════════════════

\subsection{Primary Success Criterion}

\begin{mdframed}[
  backgroundcolor=notebox,
  linecolor=headergreen,
  linewidth=1pt,
  innertopmargin=0.5em,
  innerbottommargin=0.5em,
  innerleftmargin=0.8em,
  innerrightmargin=0.8em
]
\textbf{P7-SUCCESS requires ALL of the following
at the 12-month timepoint:}

\medskip
\textbf{Criterion 1 (Adrenal):}
Mean normalised adrenal medulla mass in
Condition A $\leq$ 80\% of mean normalised
adrenal medulla mass in Condition B.
(AMR $\leq 0.80$, $p < 0.05$, two-tailed
Mann-Whitney U or independent samples
$t$-test depending on normality.)

\medskip
\textbf{Criterion 2 (HPA):}
Mean cortisol stress response ($\Delta$cortisol,
T$+$30 minus baseline) in Condition A $\leq$
75\% of mean cortisol stress response in
Condition B.
(CRR $\leq 0.75$, $p < 0.05$.)

\medskip
\textbf{Criterion 3 (Genotype confirmation):}
No significant difference in allele frequency
distribution between Condition A and Condition B
animals at any locus in the genotyping panel.
(Confirmed by STRUCTURE analysis prior to
primary endpoint analysis.)
\end{mdframed}

\subsection{Partial Confirmation Criteria}

If Criterion 1 is met but not Criterion 2,
or vice versa, the result is recorded as
\textbf{partial confirmation} with the
following interpretation:

\textbf{Criterion 1 met, Criterion 2 not met:}
Structural NCC arrest occurs (adrenal
morphology shifted) but the HPA axis
set point requires longer field exposure
or earlier (intrauterine) field application
to shift. This directs to P8 (maternal field
window experiment) as the necessary
follow-on test.

\textbf{Criterion 2 met, Criterion 1 not met:}
HPA axis reactivity shifted (functional
endpoint) without structural adrenal change.
Suggests functional reprogramming precedes
structural morphology change. Consistent
with the developmental plasticity sequence
prediction: functional endpoints respond
faster than structural ones.
Both partial outcomes are informative.
Neither is a falsification of the mechanism.

% ══════════════════════════════════════════════════════════════
\section{Control Conditions and
Confound Management}
\label{sec:controls}
% ══════════════════════════════════════════════════════════════

\subsection{The Critical Controls}

\begin{mdframed}[
  backgroundcolor=controlbox,
  linecolor=headergreen,
  linewidth=0.8pt,
  innertopmargin=0.5em,
  innerbottommargin=0.5em,
  innerleftmargin=0.8em,
  innerrightmargin=0.8em
]
\textbf{C1 --- Genotype matching (primary confound
control):}
The genetic model predicts no morphological
difference between genotype-matched animals in
different fields. Genotype matching is therefore
not merely good experimental practice --- it is
the specific instrument that rules out the
genetic model if the primary prediction confirms.
Genotype confirmation must precede any primary
endpoint analysis. If genotype matching fails,
the experiment must be repeated with a larger
cohort before any conclusion is drawn.

\medskip
\textbf{C2 --- Domestic-genome reference (Cond.\ C):}
Provides the domestic-field domestic-genome
reference point. If Condition A approaches
Condition C values, this confirms that the
field is producing the same morphological
outcome regardless of genome origin.
If Condition A does not approach Condition C
values, the field effect exists but is
incomplete for postnatal-only field application
--- which directs to P8.

\medskip
\textbf{C3 --- Litter cohort replication:}
Minimum 3 independent litter cohorts required
before primary analysis is performed.
Litter of origin is included as a random effect
in all mixed-effects models. This controls for
maternal effects within wild-type litters that
are not attributable to field assignment.

\medskip
\textbf{C4 --- Handler blinding:}
All morphological measurements, blood draws,
and sample processing must be performed by
handlers blinded to condition assignment.
Genotyping laboratory must be blinded.
Image analysis must be blinded.
Blinding is documented and confirmed
before any analysis.

\medskip
\textbf{C5 --- Nutrition equivalence:}
Condition A and B animals must receive
nutritionally equivalent diets from the
point of weaning (21 days). Differences
in ad libitum provision are a field component
(resource security is part of $N_{\text{domestic}}$)
but differences in total nutritional intake
could confound body composition endpoints.
Total caloric intake must be monitored and
reported. Body weight at each timepoint is
a covariate in all models.
\end{mdframed}

\subsection{Pre-specified Confound Interpretations}

\textbf{Confound: Litter size differences.}
Wild litters are typically larger (6--12 piglets)
than domestic litters selected for fostering
(4--8 piglets from domestic sow). Per-piglet
nutritional access may differ. Control: weigh
animals at every timepoint; include body weight
as covariate.

\textbf{Confound: Foster rejection.}
Domestic sow may reject wild-genome piglets
due to odour differences. Protocol: introduce
piglets within 6 hours of birth during foster
sow's own parturition or within 48 hours of
a litter she has accepted. Document all
rejection events. Replace foster sow if
rejection persists $> 24$ hours. Rejection
rate reported in results.

\textbf{Confound: Epigenetic carryover from
wild mother.}
The F1 animals' epigenome was established in
utero under $N_{\text{wild}}$ (wild mother).
This is not a confound --- it is the experimental
condition. The prediction is that postnatal
$N_{\text{domestic}}$ will shift morphology
despite the intrauterine wild field. If the
effect is absent, it directs to P8 (intrauterine
field required). If the effect is present, it
demonstrates that postnatal field is sufficient
for at least partial arrest. The intrauterine
starting condition is identical across Condition
A and B animals and is not a confound for the
between-condition comparison.

% ══════════════════════════════════════════════════════════════
\section{Failure Diagnosis}
\label{sec:failure}
% ══════════════════════════════════════════════════════════════

\begin{mdframed}[
  backgroundcolor=failbox,
  linecolor=red!50!black,
  linewidth=0.8pt,
  innertopmargin=0.5em,
  innerbottommargin=0.5em,
  innerleftmargin=0.8em,
  innerrightmargin=0.8em
]
\textbf{If P7 fails (no morphological difference
between Condition A and Condition B with
confirmed genotype matching), the following
diagnoses are pre-specified:}

\medskip
\textbf{Diagnosis F1 (Critical window missed):}
The NCC migration critical window is prenatal,
not postnatal. Postnatal field application
cannot arrest an already-completed developmental
program.
\textit{Direction:} Proceed immediately to P8
(maternal field from conception experiment).
P7 failure under this diagnosis does not
falsify the field-arrest model. It specifies
the critical window more precisely.

\medskip
\textbf{Diagnosis F2 (Field specification
incomplete):}
The domestic field as applied (indoor housing,
ad libitum food, human contact) is not the
complete field specification required to produce
the arrest. A missing field component (e.g.,
specific microbiome, specific tactile stimulation,
specific olfactory environment from domestic
conspecifics from conception) is required.
\textit{Direction:} Compare field component
specification between Condition A and known
domestication events. Identify the missing
component. Redesign with complete field.

\medskip
\textbf{Diagnosis F3 (Genetic floor effect):}
The wild-genome animals have a deeper Basin A
than the field-arrest model's current
specification predicts. The field applied is
sufficient to arrest a less deep Basin A
(as in heritage domestic breeds) but not
sufficient to arrest the wild-caught
\textit{Sus scrofa} Basin A within one
postnatal generation.
\textit{Direction:} Repeat with heritage
breed wild-type animals (Mangalitsa, Ossabaw
Island) which have shallower Basin A due to
prior partial domestication selection. A
gradient of effect across wild-genome depth
is a prediction of the model (see P10).

\medskip
\textbf{Diagnosis F4 (Genuine falsification):}
If F1, F2, and F3 are all ruled out by
appropriate follow-on experiments, and no
morphological shift occurs under any condition
of field specification, the field-arrest model
is falsified for the domestication syndrome.
The genetic model is the correct explanation.
This outcome would be published in full.
\end{mdframed}

% ══════════════════════════════════════════════════════════════
\section{Scoring Schedule and Timeline}
\label{sec:timeline}
% ══════════════════════════════════════════════════════════════

\begin{center}
\small
\begin{tabular}{p{2.2cm}p{9.5cm}}
\toprule
\textbf{Timepoint} & \textbf{Measurements} \\
\midrule
Birth (T$=$0) &
Genotyping (blood). Body weight.
Foster assignment documented.
Photograph (dorsal, lateral bilateral). \\[0.3em]
3 months &
Body weight. Hair coarseness score (blinded).
Ear position angle (photograph).
Cortisol stress response (blood draw protocol).
Coat pigmentation photograph. \\[0.3em]
6 months &
Body weight. Hair coarseness score.
Ear position angle. Snout morphology index.
Cortisol stress response.
Coat pigmentation. Backfat ultrasound (P2). \\[0.3em]
12 months &
Body weight. All secondary endpoints repeated.
Cortisol stress response (primary endpoint).
\textbf{Necropsy:} bilateral adrenal glands
dissected, weighed, fixed (4\% PFA).
Adrenal medulla/cortex separated and weighed
post-fixation (primary endpoint).
Histology: H\&E on adrenal sections.
IHC: chromogranin A (medulla marker),
SF-1 (cortex marker). \\
\bottomrule
\end{tabular}
\end{center}

\vspace{0.5em}
\textbf{Minimum timeline from litter birth
to primary endpoint:} 12 months.

\textbf{Minimum cohort timeline including
3-cohort replication:} 18--24 months
(cohorts staggered, not sequential).

\textbf{Estimated cost per cohort:}
Animal acquisition and housing (12 months):
\$8,000--\$15,000 per condition.
Genotyping panel: \$2,000--\$4,000.
ELISA consumables: \$1,500--\$3,000.
Histology and IHC: \$3,000--\$6,000.
Total per cohort: \$25,000--\$45,000 (3 conditions).
Total for 3-cohort replication:
\$75,000--\$135,000.

% ══════════════════════════════════════════════════════════════
\section{Statistical Analysis Plan}
\label{sec:stats}
% ════════��═════════════════════════════════════════════════════

All analyses are pre-specified. No unplanned
analyses will be reported as confirmatory.
Exploratory analyses will be clearly labelled.

\textbf{Primary analysis (adrenal mass, HPA
reactivity):}
Linear mixed-effects model with condition
(A, B, C) as fixed effect, litter cohort as
random effect, body weight at 12 months
as covariate. Primary comparison:
Condition A vs.\ Condition B.

\textbf{Normality:}
Shapiro-Wilk test on residuals.
If non-normal: Mann-Whitney U test for
primary comparison; Kruskal-Wallis for
three-condition comparison.

\textbf{Effect size:}
Cohen's $d$ (parametric) or rank-biserial
correlation $r$ (non-parametric) reported
for all primary comparisons.

\textbf{Primary significance threshold:}
$\alpha = 0.05$, two-tailed.
Bonferroni correction applied for multiple
primary endpoints (AMR and CRR):
adjusted $\alpha = 0.025$ per endpoint.

\textbf{Secondary endpoints:}
All secondary endpoints analysed identically
but reported as exploratory. No success/failure
determination made from secondary endpoints
alone.

\textbf{Longitudinal analysis:}
Repeated measures for 3-month, 6-month,
and 12-month cortisol data. Mixed-effects
model with time as a within-subject factor.
Interaction term (condition $\times$ time)
reported to assess whether field effect
increases over the measurement period.

\textbf{Pre-analysis lock:}
This statistical analysis plan is locked at
pre-registration. Any deviation from it
in the final publication must be declared
in the methods section with justification.

% ══════════════════════════════════════════════════════════════
\section{Ethical Compliance}
\label{sec:ethics}
% ══════════════════════════════════════════════════════════════

This experiment involves vertebrate animals.
The following ethical requirements apply:

\begin{itemize}[itemsep=0.15em]
  \item Institutional Animal Care and Use
    Committee (IACUC) or equivalent national
    ethics body approval required before
    any animal work begins.
  \item All animals must meet or exceed the
    minimum welfare standards of the applicable
    national legislation throughout the experiment.
  \item Condition B animals (wild-type field)
    must receive supplemental nutrition sufficient
    to maintain body weight within 10\% of
    Condition A animals. The field distinction
    is in security, not starvation.
  \item Euthanasia at 12 months by trained
    personnel using approved methods (captive
    bolt or barbiturate overdose per institutional
    protocol).
  \item Foster rejection events require immediate
    supplemental feeding within 2 hours and
    veterinary assessment. Rejected piglets must
    not be left without nutrition for $>$ 4 hours.
  \item Wild-caught source animals must be legally
    obtained with appropriate wildlife licences
    for the jurisdiction of collection.
\end{itemize}

No genetic modification is performed on any
animal in this experiment. The experiment
uses only naturally occurring genomes and
environmental field manipulation.

% ══════════════════════════════════════════════════════════════
\section{Linked Pre-Registrations
and Experimental Series}
\label{sec:linked}
% ══════════════════════════════════════════════════════════════

This pre-registration is P7 in the OrganismCore
experimental series. Linked records:

\begin{center}
\small
\begin{tabular}{p{1.5cm}p{5.5cm}p{5.0cm}}
\toprule
\textbf{ID} & \textbf{Title} & \textbf{DOI / Status} \\
\midrule
P-INVERSION-1 &
Plant Architecture Inversion via Coherence
Gradient Field Specification &
\href{https://doi.org/10.5281/zenodo.19040399}
{\texttt{10.5281/zenodo.19040399}} \\[0.3em]
LECA-A &
LECA Resurrection Fungal Arm &
\href{https://doi.org/10.5281/zenodo.18986790}
{\texttt{10.5281/zenodo.18986790}} \\[0.3em]
P7 &
Cross-Fostering / NCC Arrest
(this document) &
[Zenodo DOI pending] \\[0.3em]
P8 &
Maternal Field Window Dependence &
[Zenodo DOI pending] \\[0.3em]
P9 &
H3K27me3 at NCC Loci: ChIP-seq
Pre-Registration &
[Zenodo DOI pending] \\[0.3em]
P10 &
Breed Domestication Depth vs.\
Feralization Rate &
[Zenodo DOI pending] \\[0.3em]
P11 &
Tazemetostat in Pregnant Sow:
F1 Morphology Shift &
[Zenodo DOI pending] \\
\bottomrule
\end{tabular}
\end{center}

\vspace{0.3em}
The theory paper from which P7 was derived:

Lawson, E.R. (2026).
\textit{Domestication as Field-Specified
Developmental Arrest: A Unified Causal Mechanism
for the Domestication Syndrome Derived from
Attractor Geometry.}
OrganismCore / Zenodo.
[DOI pending]

% ══════════════════════════════════════════════════════════════
\section{Reporting Standard}
\label{sec:reporting}
% ══════════════════════════════════════════════════════════════

Results of this experiment will be reported
against this pre-registration regardless of
outcome. Specifically:

\begin{enumerate}[itemsep=0.2em]
  \item If P7 succeeds (both criteria met):
    The report will state that the primary
    prediction was confirmed, that the
    genetic model's null prediction was
    rejected, and that the field-arrest
    mechanism is supported. P8 and P9
    will proceed as follow-on experiments.

  \item If P7 partially confirms
    (one criterion met):
    The report will state which criterion
    was met, apply the pre-specified
    partial interpretation, and proceed
    to the directed follow-on experiment.

  \item If P7 fails (neither criterion met):
    The report will apply the failure
    diagnosis protocol (Section~7),
    state which diagnosis applies,
    and proceed to the directed follow-on
    experiment. The failure will be
    published in full.
    A failure in which F1, F2, and F3
    are subsequently also ruled out
    constitutes a falsification of the
    field-arrest model, which will be
    published as such.
\end{enumerate}

Selective reporting --- publishing only
the confirming result if multiple outcome
measures were taken --- is explicitly
prohibited. All pre-specified endpoints
reported in all publications arising from
this pre-registration.

% ═══════════════════════════════════════���══════════════════════
% DOCUMENT METADATA
% ══════════════════════════════════════════════════════════════
\vspace{1em}
\noindent\textcolor{rulecolor}{\rule{\linewidth}{0.8pt}}
\vspace{0.3em}

{\footnotesize
\textbf{Document ID:}
P7\_CROSS\_FOSTERING\_NCC\_ARREST\_
PRE\_REGISTRATION\_v1.0\\
\textbf{Date:} 2026-03-18\\
\textbf{Author:} Eric Robert Lawson /
OrganismCore\\
\textbf{Derivation record:}
\href{https://github.com/Eric-Robert-Lawson/
attractor-oncology}
{\texttt{github.com/Eric-Robert-Lawson/
attractor-oncology}}\\
\textbf{Theory paper:}
\texttt{Domestication\_As\_Developmental\_
Arrest.tex} (same repository)\\
\textbf{Zenodo DOI:} [to be inserted upon
Zenodo confirmation]\\
\textbf{Document version:} 1.0\\
\textbf{License:} CC BY 4.0\\
\textbf{ORCID:}
\href{https://orcid.org/0009-0002-0414-6544}
{0009-0002-0414-6544}\\
\textbf{Contact:}
\href{mailto:OrganismCore@proton.me}
{OrganismCore@proton.me}
}

\end{document}
